\subsection{INVERTIBLE ELEMENTS}

DEFINITION 6. Let E be a unital magma, \(\top\) its law of composition, \(e\) its identity element and \(x\) and \(x'\) two elements of E. \(x'\) is called a left inverse (resp. right inverse, resp. inverse) of \(x\) if \(x' \top x = e\) (resp. \(x \top x' = e\) , resp. \(x' \top x = x \top x' = e\) ).

An element x of E is called left invertible (resp. right invertible, resp. invertible) if it has a left inverse (resp. right inverse, resp. inverse).

A monoid all of whose elements are invertible is called a group.

Symmetric and symmetrizable are sometimes used instead of inverse and invertible. When the law on E is written additively, we generally say negative instead of inverse.

Examples. (1) An identity element is its own inverse.

\begin{enumerate}
\def\labelenumi{(\arabic{enumi})}
\setcounter{enumi}{1}
\tightlist
\item
  In the set of mappings of E into E an element \(f\) is left invertible (resp. right invertible) if \(f\) is a surjection (resp. injection). The left inverses (resp. right inverses) are then the retractions (resp. sections) associated with \(f\) (Set Theory, II, § 3, no. 8, Definition 11). For \(f\) to be invertible it is necessary and sufficient that \(f\) be a bijection. Its unique inverse is then the inverse bijection of \(f\) .
\end{enumerate}

Let E and F be two unital magmas and f a unital homomorphism of E into F. If \(x'\) is the inverse of x in E, \(f(x')\) is the inverse of \(f(x)\) in F. Hence, if x is an invertible element of E, \(f(x)\) is an invertible element of F.

In particular, if R is an equivalence relation compatible with the law on a unital magma E, the canonical image of an invertible element of E is invertible in E/R.

PROPOSITION 3. Let E be a monoid and x an element of E.

\begin{enumerate}
\def\labelenumi{(\roman{enumi})}
\item
  For \(x\) to be left (resp. right) invertible it is necessary and sufficient that the right (resp. left) translation by \(x\) be surjective.
\item
  For \(x\) to be invertible it is necessary and sufficient that it be left and right invertible. In that case \(x\) has a unique inverse, which is also its unique left (resp. right) inverse.
\end{enumerate}

If \(x'\) is a left inverse of \(x\) then (Proposition 1)

\[
\boldsymbol {\delta} _ {x} \circ \boldsymbol {\delta} _ {x ^ {\prime}} = \boldsymbol {\delta} _ {x ^ {\prime} \top x} = \boldsymbol {\delta} _ {e} = \mathrm{Id} _ {\mathrm{E}}
\]

and \(\pmb{\delta}_{x}\) is surjective. Conversely, if \(\pmb{\delta}_{x}\) is surjective, there exists an element \(x^{\prime}\) of E such that \(\pmb{\delta}_{x}(x^{\prime}) = e\) and \(x^{\prime}\) is the left inverse of \(x\) . The other assertions of (i) follow similarly.

If \(x'\) (resp. \(x''\) ) is a left (resp. right) inverse of \(x\) , then

\[
x ^ {\prime} = x ^ {\prime} \top e = x ^ {\prime} \top (x \top x ^ {\prime \prime}) = (x ^ {\prime} \top x) \top x ^ {\prime \prime} = e \top x ^ {\prime \prime} = x ^ {\prime \prime},
\]

whence (ii).

Remark. Let E be a monoid and x an element of E. If x is left invertible it is left cancellable; for, if \(x'\) is a left inverse of x, then

\[
\gamma_ {x ^ {\prime}} \circ \gamma_ {x} = \gamma_ {x ^ {\prime} \top x} = \gamma_ {e} = \mathrm{Id} _ {\mathrm{E}}
\]

and \(\pmb{\gamma}_{x}\) is injective. In particular, if \(x\) is left invertible, the left and right translations by \(x\) are bijective. Conversely, suppose that \(\pmb{\gamma}_{x}\) is bijective; there exists \(x' \in E\) such that \(xx' = \pmb{\gamma}_{x}(x') = e\) ; then \(\pmb{\gamma}_{x}(x' x) = (xx') x = x = \pmb{\gamma}_{x}(e)\) and hence \(x' x = e\) , so that \(x\) is invertible. We see similarly that if \(\delta_{x}\) is bijective, \(x\) is invertible.

PROPOSITION 4. Let E be a monoid and x and y invertible elements of E with inverses \(x'\) and \(y'\) respectively. Then \(y' \top x'\) is the inverse of \(x \top y\) .

This follows from the relation

\[
\left(y ^ {\prime} \top x ^ {\prime}\right) \top (x \top y) = y ^ {\prime} \top \left(x ^ {\prime} \top x\right) \top y = y ^ {\prime} \top y = e
\]

and the analogous calculations for \((x\top y)\top (y^{\prime}\top x^{\prime})\)

COROLLARY 1. Let E be a monoid; if each of the elements \(x_{\alpha}\) of an ordered sequence \((x_{\alpha})_{\alpha \in A}\) of elements of E has an inverse \(x_{\alpha}'\) , the composition \(\bigcap_{\alpha \in A} x_{\alpha}\) has inverse \(\bigcap_{\alpha \in A'} x_{\alpha}'\) , where \(A'\) is the totally ordered set derived from A by replacing the order on A by the opposite order.

This corollary follows from Proposition 4 by induction on the number of elements in A.

In particular, if \(x\) and \(x'\) are inverses, \(\frac{\pi}{\top} x\) and \(\frac{\pi}{\top} x'\) are inverses for every integer \(n \geqslant 0\) .

COROLLARY 2. In a monoid the set of invertible elements is stable.

PROPOSITION 5. If in a monoid \(x\) and \(x'\) are inverses and \(x\) commutes with \(y\) , so does \(x'\) .

From \(x \top y = y \top x\) , we deduce \(x' \top (x \top y) \top x' = x' \top (y \top x) \top x'\) and hence \((x' \top x) \top (y \top x') = (x' \top y) \top (x \top x')\) , that is \(y \top x' = x' \top y\) .

COROLLARY 1. Let E be a monoid, X a subset of E and X' the centralizer of X. The inverse of every invertible element of X' belongs to X'.

COROLLARY 2. In a monoid the inverse of a central invertible element is a central element.

